\section{Related work}
\label{sec:related}

Agner Fog's instruction tables \citep{fog2025instruction} are the model
for the artefact: exhaustive per-instruction latency and throughput for
x86 microarchitectures, maintained as data rather than narrative, with
the measurement method documented separately \citep{fog2025microarch}.
This work transplants the form to a GPU and adds two things the x86
tables do not carry: per-row provenance binding each number to the
benchmark commit and results line that produced it, and an explicit
deviation column against published priors.

GPU microbenchmarking begins for our purposes with Wong et
al.\ \citep{wong2010demystifying}, who established pointer-chase latency
measurement on GT200, and Mei and Chu \citep{mei2017dissecting}, whose
fine-grained pointer chase resolved cache geometry on Kepler and Maxwell.
Jia et al.\ characterised Volta \citep{jia2018volta} and then Turing via
the Tesla T4 \citep{jia2019turing}. Their instruction-latency tables are
the priors this paper compares against, row by row, where the prior
applies. The applicability boundary matters: the T4's TU104 shares the
TU102's streaming multiprocessor, so core-domain priors transfer, but its
L2, DRAM system, and lack of NVLink make memory-system and interconnect
priors non-binding, a distinction this paper records per row rather
than per paper.

Two methodological differences from Jia et al.\ are deliberate. First,
their measurements were taken with a custom SASS assembler, giving direct
control of the instruction stream. This work stays within the public
toolchain (PTX inline assembly under \texttt{nvcc} 13.3) and instead
\emph{verifies} the emitted stream, gating every benchmark on a
disassembly check. The compiler fights this approach: constant
folding, uniform-datapath conversion, strength reduction, and
approximate-arithmetic algebra each deleted or rewrote measurement
kernels during development. \Cref{sec:methodology} catalogues the
countermeasures, because any future toolchain-bound measurement effort
will face the same opponents. Second, where their report states results,
this one also pre-registers the decision-grade questions
(\cref{sec:hypotheses}) and publishes the failures.
